% 第 20 章　电势：电场中的地形图
\chapter{电势：电场中的地形图}

\step{反常识开场}

\begin{mainline}
科技馆里有一件经典展品：范德格拉夫起电机。一个大金属球，观众把手扶在球上，工作人员摇动手柄，球上的电压一路升到二十万伏。观众的长发一根根竖起来，像炸毛的刺猬——而人毫发无损，走下台时只是裤腿有点静电。

现在把场景切回家。墙上那个普通的插座，电压只有二百二十伏，是静电球的千分之一。可是从小到大，所有大人都告诫你：绝不能碰，碰了会死人。

二十万伏摸上去没事，二百二十伏却要命。如果说“伏特数越大越危险”，这两件事必有一件在撒谎；如果说“伏特数无所谓”，那电压表上那个数到底在度量什么？

再看两件怪事。第一件：一只鸟站在几十万伏的高压输电线上一动不动，怡然自得；可如果它展开翅膀时一脚踩线、一脚碰上铁塔，会当场毙命。同一根电线，同一个电压，差别在哪里？第二件：雷雨夜，云和地面之间的电压高达上亿伏，闪电为什么不笔直地走最短路线，偏要扭出一条弯弯曲曲、还带着分叉的折线？它像在打探什么。

还有一个你亲手就能做的怪事：把塑料梳子在头发上摩擦几下，靠近桌上的碎纸屑，纸屑会隔着空气跳起来，贴上梳子。梳子对纸屑做了功——把纸屑提起来了。这份功是谁付的？只能是你摩擦梳子的那几下手。可是你的能量是怎么“存”在梳子周围的空气里，等纸屑一来就兑现的？

这一章的任务，就是把“伏特”这个被用滥了的词拆开。我们会给电场画一张地图——一张看不见的地形图。画完这张图，开场的每一桩怪事都会自动解开。
\end{mainline}

\step{先猜一猜}

\begin{mainline}
在往下读之前，请先对下面三句话给出你自己的直觉判断，最好写在纸上：

\begin{itemize}
  \item 静电球上二十万伏却电不死人。你猜原因是：电压表读数是虚的？是“电的量”太少？还是另有原因？
  \item 鸟站在高压线上安然无恙。你猜关键是：鸟的脚不导电？电线外皮绝缘？还是鸟身体两端的“某种差别”几乎为零？
  \item 气球在毛衣上蹭几下就能粘在墙上，但过几个小时或几天，它会自己掉下来。你猜它是“电力用完了”，还是“贴上去的劲”消失了——这两者是一回事吗？
\end{itemize}

我提前剧透一句：这三题的答案都指向同一个概念。这个概念你已经非常熟悉——爬山时你每天都在和它打交道，只是它在电场里换了个名字。
\end{mainline}

\step{动手实验}

\begin{experiment}[看见那张“看不见的斜坡”]
你需要：一把塑料梳子（或一只气球）、一堆碎纸屑（把废纸撕成指甲盖大小）、一个能放出极细水流的水龙头。天气越干燥，效果越好。

第一步，隔空做功。把梳子在干燥的头发（或毛衣）上用力摩擦十几下，慢慢靠近纸屑，不要碰到。在某个距离上，纸屑会突然跳起，啪地贴上梳子。注意观察：纸屑是\emph{加速}冲上去的——它在空中就被推动了，不是碰到梳子才被黏住。隔着空气，梳子对纸屑做了功。

第二步，看“地形”如何弯折水流。把水龙头调成一条极细的连续水流（像棉线那么细），重新摩擦梳子，靠近水流侧面，同样不要碰到。你会看到水流明显地弯向梳子。水分子并没有被梳子“抢走电子”——整条水流完好地流着——但水里的每一个分子都被隔空推了一把。

第三步，等纸屑“翻脸”。让纸屑在梳子上贴几秒，有些纸屑会突然跳开，像被梳子嫌弃了。先吸上来，再推开去——吸引力怎么变成了排斥力？（如果你在第 19 章读过接触起电，这里已经能猜到答案：纸屑贴上后从梳子那里分到了同种电荷。）

整个实验里你最该带走的感受是：梳子周围似乎存在一张看不见的地形。纸屑和水流一进入这张地形，就往某个方向“滑”。能量是你手臂付的，存在这张地形里，等一个带电的家伙闯进来，就当场兑现。本章剩下的篇幅，就是把这张地形图画出来，并标上刻度。
\end{experiment}

\step{旧模型遇到麻烦}

\begin{mainline}
上一章（第 19 章）交到我们手里的工具是库仑定律和电场：知道了电荷的分布，就能算出空间每一点的电场，知道电场就能算出任何一个试探电荷受的力，然后一步步推它的运动。理论上，这套工具足够描述静电世界里发生的一切。那么，拿它去解决本章的问题，会发生什么？

麻烦一：算不动。电场是随位置变化的，试探电荷每挪动一小步，力的方向和大小都得重新算一遍。追踪一个电子在两块带电金属板之间的旅程，就像在没有地图的山区里赶路：你只能低头看脚底的坡度，一步一测，永远不知道前方三里地外是什么在等着你。物理学家不喜欢这种活法——我们需要的是地图，不是盲人摸象。

麻烦二：一笔奇怪的账。有人（历史上是十八世纪的数学家们）用库仑定律算了一笔账：把一个试探电荷从电场中的 A 点搬到 B 点，无论走哪条路——直的、弯的、绕地球一圈再回来的——电场力做的总功竟然完全一样。更惊人的是，沿任何一条闭合路线走一圈回到出发点，总功严格等于零。这不是巧合。你在爬山时遇到过同样的事：无论走缓坡还是攀岩，从山脚到山顶，重力对你做的功只看高度差；兜一圈回到原地，重力做功恰好为零。凡是有这种“只认终点不认路”性质的力，背后都藏着一张“高度表”。

麻烦三：能量藏在哪里。梳子把纸屑吸起来，纸屑获得了动能；可动能不会凭空出现，它一定是从某个“库存”里取出来的。库存在哪？梳子摸上去和普通梳子没有任何不同。账本上有一笔能量，寄存在梳子周围的空气里，但旧模型的语言里没有一个词能称呼它。

三个麻烦指向同一条出路：别再逐点问“这里的电荷受多大的力”，改问“这里比那里‘高’多少”。只要电场也有“高度”这种东西，地图就有了，账本就有了，能量库存也有了。问题是：电场里的“高度”到底是什么？
\end{mainline}

\step{发明新概念}

\begin{mainline}
像历史上的物理学家一样，我们一层一层把这个概念造出来。

第一层，\textbf{电势能}。把电荷放在电场中某一点，它就带着一份“潜在的能”——只要松手，电场力就会推着它跑，把这份能兑现成动能。我们规定：电场力做了多少正功，电势能就减少多少。这和你熟悉的重力势能是同一个逻辑：石头下落，重力做功，势能减少，动能增加。电势能是“电荷和电场这对搭档”共有的库存。

第二层，\textbf{电势}。电势能有个麻烦：同一个位置上，放一库仑电荷和放两库仑电荷，库存差一倍——它既依赖于位置，也依赖于你放上去的是谁。而地图应该只描述地形本身，与来爬山的是乒乓球还是铅球无关。于是我们把电势能除以电荷量，定义：电势 \(V\) 等于单位电荷在该点的电势能。电势是空间本身的属性（严格说，由产生电场的那些源电荷决定），与有没有试探电荷在场无关。单位是焦耳每库仑，这个单位有个专门的名字：伏特。

第三层，\textbf{电势差}，也就是日常说的“电压”：两点的电势之差。为什么要专门发明它？因为世界上所有能造出来的电压表，量的永远是差——两根表笔，一头一个点。绝对电势的零点可以随便定（通常约定无穷远处为零，工程上约定大地为零），但两点之间的差是铁打的物理事实。就像你可以把海拔零点定在海边或定在你家客厅地板，珠穆朗玛峰和你家房顶的高度差却不变。

现在，请出本章的核心比喻：电势是电场中的“海拔”，电势分布是电场中的“地形图”。按照全书的规矩，先说它像什么，再说它不像什么。

它像的地方：电势像海拔；电场像坡度——坡越陡，电场越强，而且电场指向“下坡”的方向；正电荷像山坡上的小球，被推向低电势处；沿等高线走不费功，沿等势面移动电荷，电场力不做功。

它不像的地方，请务必记牢。第一，海拔几乎总是正的，电势可正可负：负电荷周围是一口“井”，电势比无穷远低。第二，山坡上滚下来的是真实的小球，电场里却没有任何“电势”在流动——电压本身不流动，流动的是电荷；说“电压流过去了”就像说“高度差流过去了”一样不通。第三，负电荷是一颗“反重力小球”：正电荷被推向低电势，负电荷被推向高电势——同一张地形图，两种走法相反。第四，山坡是二维曲面，电势地形却充满整个三维空间，而且完全看不见，只能靠试探电荷去“踩点”。

在往下走之前，不妨把这套新概念对进全书反复追问的三个问题里。系统此刻是什么状态？答：空间每一点有一个电势值，整张地形图就是静电系统的状态。什么规律决定状态如何变化？答：电荷的分布决定地形，而地形一旦给定，任何电荷受的力都由它所在处的坡度决定。我们怎样通过测量知道它？答：电压表——它有两根表笔，分别“踩”在两个点上，报出的读数就是那两点的落差。地图看不见，但落差是可以一格一格量出来的。

有了地形图，就能画出它的“等高线”了：把电势相等的点连起来，得到\textbf{等势面}（平面上是等势线）。沿等势面走，电场力永远不做功——这说明电场力在等势面方向上没有分量，也就是说：\textbf{电场处处与等势面垂直}，正如山坡的最陡方向处处与等高线垂直。电场线总是从电势高的地方指向电势低的地方（这是对正试探电荷而言的“下坡”），而且等势面越密集的地方，坡越陡，电场越强。
\end{mainline}

\begin{figure}[htbp]
\centering
\begin{tikzpicture}[scale=1.1]
  \fill (0,0) circle (2.5pt) node[below=4pt] {$+Q$};
  \foreach \r in {0.9,1.7,2.5}
    \draw[dashed] (0,0) circle (\r);
  \foreach \a in {0,45,90,135,180,225,270,315}
    \draw[-{Stealth},thick] (\a:0.25) -- (\a:2.9);
  \node at (2.05,1.15) {等势面（虚线）};
  \node at (2.3,-1.9) {电场线（实线）};
\end{tikzpicture}
\caption{正点电荷周围的地形图：虚线是一圈圈等势面，像山峰的等高线，越往外越“矮”；实线电场线沿径向指向外，处处与等势面垂直——那就是“下坡”最陡的方向。}
\end{figure}

\step{一句话模型}

\begin{oneline}
电势是电场的海拔：电场力搬运电荷做的功只认起点与终点的高度差、不认路，这个高度差就是电压；正电荷被推向低处，负电荷被推向高处，而“储存起来的高度差”就是电容器里的能量。
\end{oneline}

\step{数学翻译器}

\begin{mainline}
先把三句直觉翻译成三个式子。每一个式子，我们都按老规矩问四个问题：它描述什么？为什么这样写？何时成立？何时失效？

第一句：\textbf{电场力做的功} \(=\) \textbf{多少电荷} \(\times\) \textbf{落差}。写成
\[
W_{A\to B} = q\,(V_A - V_B) .
\]
它描述：电场力把电荷 \(q\) 从 A 点搬到 B 点做了多少功。为什么这样写：账就该这么记——功是“多少货”乘“每单位货的高度差”，\(q\) 是货，\(V_A - V_B\) 是落差；这也正是一开始定义电势时想要的性质。何时成立：静电场中，即电场不随时间变化、也没有变化的磁场掺和。何时失效：本章“模型的边界”一节会揭晓，这个式子脚下的地基有一条裂缝。

立刻做特殊情形检查。取 \(q\) 为负：右边自动变号——负电荷的“下坡”方向与正电荷相反，式子不用改，符号自己翻转，通过。取 \(V_A = V_B\)（两点在同一等势面上）：\(W=0\)，无论中间绕了多远的路，通过——这正是“等高线上不费劲”。反过来，如果测得 \(W=0\) 而 \(q\neq 0\)，则必有 \(V_A = V_B\)：不费劲的路一定没爬坡，通过。

第二句：\textbf{电场就是电势的坡度}。在匀强电场里（比如平行板电容器两板之间），沿电场方向相距 \(d\) 的两点，电势差为
\[
V = E\,d .
\]
它描述：电势沿电场方向“每走一米降多少伏”，这个速率就是场强 \(E\)。所以 \(E\) 的单位既可以写牛每库仑，也可以写伏每米——“单位电荷受多大的力”和“电势降得多快”是同一件事的两种说法。为什么这样写：坡度等于落差除以水平距离，这是任何地图的常识。何时成立：匀强电场；电场不均匀时，这个式子要对“很小一段距离”成立，那是数学桥层的事。何时失效：随时间剧烈变化的场里要小心，原因同上。

检查：\(d=0\) 时 \(V=0\)——原地没有落差，通过。\(d\) 加倍 \(V\) 加倍——走两倍的路，爬两倍的坡，通过。方向反转（沿电场反方向走），\(V\) 变号——从下坡变上坡，通过。还可以把这个关系画成一张真正的“剖面图”：以位置为横轴、电势为纵轴，平行板之间是一条笔直下降的斜线，斜线的陡度就是场强。看地形图的人扫一眼坡度就知道风往哪吹，看这张剖面图的人扫一眼斜率就知道电场有多强、朝哪边。
\end{mainline}

\begin{figure}[htbp]
\centering
\begin{tikzpicture}[scale=1.2]
  \draw[-{Stealth}] (0,0) -- (4.6,0) node[right] {位置 \(x\)};
  \draw[-{Stealth}] (0,0) -- (0,2.9) node[above] {电势 \(V\)};
  \draw[thick] (0,2.4) node[left] {\(V_{+}\)} -- (3.8,0.3);
  \draw[dashed] (3.8,0.3) -- (3.8,0) node[below] {\(d\)};
  \draw[dashed] (3.8,0.3) -- (0,0.3) node[left] {\(V_{-}\)};
  \node at (2.1,1.75) {斜线的陡度 \(=\) 场强 \(E\)};
\end{tikzpicture}
\caption{平行板之间的“电势剖面图”：电势沿电场方向均匀下降，是一条直线。直线越陡，同样的距离内落差越大，电场越强——电场就是这张图的坡度。}
\end{figure}

\begin{mainline}
来个真实数字找感觉：干燥空气的击穿场强约 \(3\times10^{6}\) 伏每米。两根相距一毫米的导线，电压差超过约三千伏就可能打出火花；脱毛衣时的电火花有几毫米长，说明电压上万伏——但你安然无恙，原因到“回头解释开场现象”再算这笔账。
\end{mainline}

\begin{figure}[htbp]
\centering
\begin{tikzpicture}[scale=1.15]
  \draw[thick] (-2,1.5) -- (2,1.5) node[right] {$+Q$ 板};
  \draw[thick] (-2,-1.5) -- (2,-1.5) node[right] {$-Q$ 板};
  \foreach \x in {-1.5,-0.9,...,1.5}
    \draw[-{Stealth},thick] (\x,1.35) -- (\x,-1.35);
  \foreach \y in {-0.75,0,0.75}
    \draw[dashed] (-1.9,\y) -- (1.9,\y);
  \node at (-2.55,0) {虚线：等势线};
\end{tikzpicture}
\caption{平行板之间是匀强电场：实线电场线笔直、均匀、竖直向下，虚线等势线是水平等间距的平行线——这是一张“均匀斜坡”的地图，越往下电势越低。}
\end{figure}

\begin{mainline}
第三句：\textbf{点电荷的电势}。以无穷远处为零点，点电荷 \(Q\) 在距离 \(r\) 处制造的电势，距离越远越低，而且“衰减速率”比力慢——力按 \(1/r^2\) 衰减，电势只按 \(1/r\) 衰减：距离加倍，力减到四分之一，电势只减到一半。精确的式子需要一点乘法以外的工具，放进数学桥层；但你可以先记住图像：正电荷是一座无限高的尖峰，负电荷是一口无限深的井，等高线都是同心圆，靠得越近越密——坡越陡。
\end{mainline}

\begin{mathbridge}
点电荷电势的精确式为
\[
V(r) = \frac{kQ}{r},\qquad k = 9\times10^{9}\ \text{牛·米}^2/\text{库}^2 .
\]
它和场强公式 \(E(r) = kQ/r^2\) 的关系，正是“高度”与“坡度”的关系：电势是场强沿径向一路累积起来的（积分），场强是电势沿径向的变化率（负的导数）。这就是为什么电势比力衰减得慢一次幂。

用它回答四个问题。描述什么：点电荷周围地形的高度表。为什么这样写：把库仑力沿径向从无穷远累积到 \(r\)，每段路上的功加总，恰好得到 \(kQ/r\)。何时成立：点电荷、或球对称分布的电荷在其外部（这和引力的壳层结论如出一辙）。何时失效：看极限——\(r\to\infty\) 时 \(V\to 0\)，与零点约定自洽，很好；但 \(r\to 0\) 时 \(V\to\infty\)，这是模型在拉警报：任何真实电荷都有大小，当 \(r\) 小于电荷本身的尺寸，“点电荷”假设先垮掉，公式随之作废。公式的发散不是世界的灾难，而是模型的边界在自己说话。
\end{mathbridge}

\begin{mainline}
第四句：\textbf{电容器存的是“分开的状态”，不是电荷}。两块平行金属板，一块带 \(+Q\)、一块带 \(-Q\)，板间电场近似匀强。定义电容
\[
C = \frac{Q}{V},
\]
即“每升高一伏电压能装下多少电荷”。它描述：这个装置储存“电荷分离”的胃口有多大。为什么这样写：电压正比于电荷是实验事实（在材料不变时），比例系数就是电容。何时成立：电介质不被击穿、材料性质不变。何时失效：电压高到击穿绝缘层，或电荷密度高到材料吃不消。

存的能量是多少？结论先行：
\[
U = \frac{1}{2}QV = \frac{1}{2}CV^2 .
\]
那个 \(\frac{1}{2}\) 从哪来？想象你一批一批地往板上搬电荷。搬第一批时，板间还没有电压，几乎不用费劲；搬到一半时，电压已有 \(V/2\)；搬最后一批时，要顶着满满的电压 \(V\) 硬推上去。每一批电荷经历的“平均落差”只有 \(V/2\)，所以总功是平均电压乘总电荷：\(\frac{1}{2}QV\)。检查特殊情形：\(V=0\) 时 \(U=0\)，没充电当然没储能，通过；\(V\) 加倍时 \(U\) 变成四倍而不是两倍——既多装了一倍的货，又把每一批货抬到了平均两倍的高度，两个因素叠乘，通过。

这笔能量存在哪里？不在金属板里，不在导线上，而在\textbf{两板之间的电场里}。请注意一个常被说错的点：电容器并没有“存电荷”——两块板的净电荷之和仍是零，正等于负。它存的是“正负电荷被硬生生分开”这个状态，以及把它们拉开时你（或电池）付出的能量。这和本书反复强调的一句话是同源的：不说“电被电池用完了”，能量在转移和寄存，电荷本身始终守恒。

真实数据估算一：相机闪光灯。典型闪光灯电容约 \(150\) 微法，充电电压约 \(300\) 伏：
\[
U=\frac12\times 1.5\times10^{-4}\times 300^2 \approx 7\ \text{焦} .
\]
七焦耳是什么概念？大约是一个苹果从七米高处落地时的动能。这七焦耳在约千分之一秒内全部倒进灯管——所以闪光灯才能又亮又快。它的亲戚、医院里的除颤器，一次电击释放一百五十到三百六十焦，是把几十倍于闪光灯的能量在几毫秒内推过胸腔，让紊乱的心脏电活动“重启”。同一种地形图，一个救急，一个救命。

真实数据估算二（也是本章留给后面的伏笔）：把一个电子（电荷 \(e = 1.6\times10^{-19}\) 库）从电势低一伏的地方搬到高一伏的地方，需要 \(1.6\times10^{-19}\) 焦。这份能量小到在力学里没有名字，但在原子里却是标准零钱——科学家干脆把它当单位用，叫\textbf{电子伏特}。化学键的成与断、电池的一格一格电压、原子发光的颜色，全都是几个电子伏特量级的事。后面讲电池（第 21 章）、讲原子发光时会天天用到它。
\end{mainline}

\begin{mainline}
最后补上微观的一块拼图：\textbf{电介质与极化}。把一张纸、一块玻璃插进电容器两板之间，电容会变大——同样的电压能装下更多电荷。为什么？

绝缘体里的电子被各自的原子核拴着，不能像导体里那样长跑。但电场仍然能“推一把”：每个分子的正电荷中心被推向低电势侧，负电荷中心被推向高电势侧，两者错开一丁点距离——整个分子变成一头正一头负的小偶极子，这个过程叫\textbf{极化}。有些分子（比如水分子）天生就一头正一头负，电场做的则是让它们集体转向、整齐列队。动手实验里水流被梳子吸弯，正是水分子被极化、转向的证据。

亿万个小偶极子整齐排列的效果是：在电介质贴着金属板的表面上，浮现出一层“极化电荷”，极性与板上电荷相反。它们部分抵消了板上电荷制造的电场——板间坡度变缓了，同样的 \(Q\) 对应更小的 \(V\)，于是 \(C=Q/V\) 变大。变大的倍数叫介电常数：纸约二到三倍，玻璃约五到十倍，水（静电情形）约八十倍，某些陶瓷可达数千倍。电容器能做得那么小，全靠电介质帮忙。

极化也有极限。电场太强时，电子会被直接从分子里扯出来，绝缘体瞬间变成导体——这就是\textbf{击穿}。空气的击穿场强前面用过，约 \(3\times10^{6}\) 伏每米。所以每只电容器外壳上都标着耐压值：那不是建议，是地形图的海拔上限。

电介质就在你身边默默工作。手机屏幕底下那层玻璃就是电介质：你的手指并不真正“碰到”电路，只是靠近时，手指里的电荷让玻璃极化，局部电容发生一丁点变化，芯片读出这个变化，就知道你点了哪里。冬天脱毛衣的火花、梳子和水流、触摸屏、相机闪光灯，看似四件不相干的事，在同一张地形图上是四个相邻的地标。
\end{mainline}

\begin{challenge}
学过微积分的读者可以把“坡度”说得更彻底。沿 \(x\) 方向，场强是电势的负导数：
\[
E_x = -\frac{dV}{dx}.
\]
在三维里，电场是电势的负梯度：\(\mathbf{E} = -\nabla V\)。负号的意思是“电场指向下坡”。一旦接受这个语言，本章所有定性结论都变成一行行恒等式：等势面上 \(dV=0\)，故 \(\mathbf{E}\) 必与其垂直；“绕一圈不做功”对应 \(\oint \mathbf{E}\cdot d\mathbf{l}=0\)，这是静电场版的“地形图可画性条件”。而当这个环路积分可以不为零时（第 23 章，变化的磁场登场），电势这张地图在数学上就画不出来了——不是物理学家不想画，是那块“地形”拧成了沿闭合回路持续下坡的怪圈，像埃舍尔画里那座永远走不到顶的楼梯。真实世界里确实存在这样的场，到时候我们换一套语言。
\end{challenge}

\step{模型的边界}

\begin{mainline}
每张地图都有出版范围，电势这张地形图有四条边界。

边界一：电势的零点可以随便定，只有电势差是物理的。说“这一点电势一千伏”却不说明相对谁，等于报海拔不报零点。工程上统一取大地为零点，这就是“接地”的来历——把设备外壳和大地用导线连起来，强行让它和所有人脚下的“海拔零点”等势。这也是为什么相对的两根输电线之间差几十万伏、而站在地上的人只要不同时碰到它们就没事。

边界二：“电压是电流的压力”——这个比喻用在哪里会翻车。它有用的地方：电势差确实会驱动电荷定向移动，像水压驱动水流；几节电池串联像把几层楼的高度差叠起来。但请注意它不像的地方，一共四条。其一，电压不是力：它的单位是焦每库，是一记账目，不是一牛顿推力。其二，水压需要水管里有水才能存在，电势差却可以横跨真空——平行板之间空无一物，电压照样在。其三，断开电路时，电池两端的电压依然存在，压力比喻很难想象“关死阀门还能一直维持压差”，但高度差比喻毫无障碍：一楼和二楼的高度差，与楼梯上有没有人在走毫无关系。其四，有了电压并不等于就有了多大的电流——那还取决于“路况”（电阻），这是下一章的主角。电压只是落差，不是流量。

边界三：电势存在的地基是“绕闭合回路一周，电场力不做功”。这个地基在静电场里坚如磐石，但它不是宇宙定律。第 23 章我们会亲眼看到，当磁场随时间变化时，会产生一种涡旋状的电场，沿闭合回路走一圈，电场力做的功可以不为零——在那种场里，“电势”连定义都立不住，只能退回电场本身说话。这不是本章模型的瑕疵，而是它的适用范围。一张地图画不出整个地球，不等于地图错了。

边界四：一组专门治疗直觉的反例——\textbf{电场为零的地方，电势可以很高；电势为零的地方，电场可以很强}。两个等量同号正电荷连线的中点，两边的推力恰好抵消，电场为零；但电势是两份正值相加，比无穷远高得多——把一个正电荷从无穷远推到中点，一路上坡，只是终点恰好是一块平台。反过来，等量异号电荷连线的中垂面上，电势处处为零（正峰的“高”和负井的“深”恰好抵消），但电场处处不为零——那是一片海拔为零、坡度却很陡的平原。高原上海拔三千米的平地，和海边陡崖，把这件事说得再明白不过：“高度”和“坡度”是两笔账。
\end{mainline}

\step{回头解释开场现象}

\begin{mainline}
现在带着地形图，回到本章开头的四个悬念。

悬念一：为什么二十万伏的静电球电不死人，二百二十伏的插座却要命？先算静电球的账。金属球半径约二十厘米，电容只有约二十皮法（\(2\times10^{-11}\) 法），充到二十万伏时，球上的总电荷 \(Q=CV\approx 4\times10^{-6}\) 库——只有几微库；储能 \(\frac12 CV^2\approx 0.4\) 焦，而且一次放完，没有后续。这几微库电荷流过你身体进入大地，只够让你皮肤一麻、头发竖起。再看市电：二百二十伏的“落差”确实矮得多，但插座背后连着整个电网——只要你保持接触，电荷可以源源不断地流，每秒几焦、几十焦地持续往你身体里灌。决定危险程度的从来不是电压这一个数，而是“落差乘以总量、再乘以持续时间”这笔完整的能量账。伏特告诉你每一库仑会经历多大的落差，却不告诉你总共会来多少库仑。

悬念二：鸟站在高压线上为什么没事？因为它的两只脚踩在同一根导线上，两点之间几乎没有电势差——没有落差，就没有“水流”经过身体。触电的本质从来不是“碰到了高电压”，而是“身体两端出现了电势差”：一脚踩线、一脚碰塔，等于把自己接在了落差几十万伏的两点之间。高压线上的带电作业工人穿的等电位服，就是把全身用导电纤维连成一体，强制全身等势——让身体变成等势面的一部分，电流就没有理由穿过人体。

悬念三：闪电为什么走折线？云和地面像一对拉得极远的平行板，但算一下就知道：电势差 \(10^{8}\) 伏除以两三公里的距离，平均场强只有约 \(5\times10^{4}\) 伏每米，远低于空气的击穿场强。所以整块空气不会同时击穿，闪电只能“摸着走”：某处的水滴、冰晶或离子先撕开一小段导电通道，云的“高电势”就顺着通道被带到了离地面更近的地方——电势地形被重新画了一遍；新的最陡处再撕开下一段，地形再改一次。击穿通道和电势地形互相改写、一步一步向下摸，走出来的自然是弯弯曲曲、带着分叉的路线。折线不是闪电“选”出来的，是地形逐段塌陷的遗迹。

悬念四：气球粘墙几天为什么会掉？气球靠极化贴在墙上：它的电场让墙面分子极化，异性电荷相吸，正如梳子吸弯水流。它掉下来不是因为“电被用完了”——这个说法在本书里是禁句——而是气球表面的电荷被空气里的离子和潮气一点点中和、被灰尘一点点带走。电荷没了，极化没了，吸引力没了，高度差抹平了，气球自然滑落。消失的是“分开的状态”，不是电本身。

四桩怪事，一张地形图全部收账。这就是新概念是否值钱的检验标准。
\end{mainline}

\step{脑洞实验与挑战题}

\begin{thought}[把雷雨云和地面当成一块平行板电容器]
一次典型的闪电：云地之间电势差约 \(10^{8}\) 到 \(10^{9}\) 伏，单次输送电荷约五到二十库仑。用地形图的第一条公式算能量：
\[
W = qV \approx 20\times 5\times10^{8} = 10^{10}\ \text{焦} \approx 2700\ \text{千瓦时},
\]
够一个三口之家用一年。听起来是座金山？可惜这笔能量摊在几公里长的通道上、几十微秒内放完，绝大部分用来把空气加热到三万度、发出巨响——真正可以收割的所剩无几。有人认真算过“闪电发电”的经济账：即便把落在一座城市的闪电全部接住，一年也不过相当于一座小电厂几天的输出，还要先解决“如何接住一亿伏”这种小问题。

再把尺度倒过来。想让一毫米的空气隙直接击穿，按击穿场强 \(3\times10^{6}\) 伏每米算，需要约三千伏——手机充电器插脚插入插座时偶尔打的小火花，正是这个量级。从微米级的芯片绝缘层，到公里级的雷雨云，\(V=Ed\) 这一条式子跨了九个数量级依然管用。一张地图能用的尺度跨度，就是它最好的成色证明。
\end{thought}

\begin{puzzles}
\item 两个等量同号的正电荷，连线中点处电场恰好为零。该点的电势也是零吗？把一个正试探电荷从无穷远处沿连线慢慢推到中点，电场力对它做正功还是负功？（提示：电场为零只说明中点“此刻不推人”，不代表一路上没人推。想象一段上坡的终点恰好是一块平台——平台本身是平的，但它仍然比山脚高。电势记的是全程的账。）
\item 避雷针为什么做成尖的，而不是装一个光滑的圆球顶？（提示：导体表面整体是一个等势面；尖端会把附近的等势面“挤”得弯曲而密集——同样的落差被塞进更短的距离，坡度陡得多，空气最先在那里被击穿。闪电于是被“请”到针上，再沿导线引入大地，而不是随机挑一处屋顶。）
\item 电容器的储能是 \(\frac12 QV\)，不是 \(QV\)。假如改成“先把板间电压充到 \(V\)，再把全部电荷 \(Q\) 一次性放上去”，账不就变成 \(QV\) 了吗？现实中为什么做不到？（提示：电荷是一批一批搬上去的，越往后板间电压越高、搬运越费劲，平均落差只有 \(V/2\)。而“先有电压 \(V\)”要求板上已经有电荷——可那时你还没开始搬。先有鸡还是先有蛋，账本上自有答案。）
\item 平行板电容器充好电后与电池断开（电荷无处可逃），然后把两板间距拉大。板间电场、电压、储能分别怎么变？如果储能变了，差额是谁付的？（提示：\(Q\) 不变；匀强电场的强弱只由板上的电荷密度决定，与板距无关；而 \(V=Ed\)。储能若增加，找找看你的手在拉板子时要对抗什么。）
\end{puzzles}
